\section{自动驾驶商业模式：谁拥有数据和客户}

讲完 Waymo 和 Momenta 的个人训练后，需要抽象到商业模式层面。因为自动驾驶和机器人一样，技术路线从来不是纯技术选择，而是和数据入口、客户关系、利润池绑定在一起。

高继扬在访谈中把自动驾驶商业模式分成几类：Waymo 式 robotaxi，车企的卖车加软件订阅，Momenta 式供应商，以及类似华为的整车利润平台。这个分类帮助我们理解技术路线为什么会分叉：谁拥有车、谁拥有用户、谁拥有数据、谁承担服务责任，会直接影响技术架构和组织能力。

\begin{figure}[H]
\centering
\includegraphics[width=0.94\textwidth]{autonomous-driving-business-models.png}
\caption{自动驾驶的四种商业模式：Robotaxi、车企订阅、供应商、整车利润平台。}
\end{figure}

\begin{knowledgebox}{读图：商业模式决定数据循环}
Robotaxi 自营车队，数据和服务都在自己手里，但规模化慢；车企订阅拥有车辆和用户，能直接拿到数据；供应商要通过车企项目进入量产，客户价值和数据回流都依赖合作；整车利润平台则通过品牌、渠道、智驾和座舱重定义整车利润池。
\end{knowledgebox}

\subsection{为什么这对机器人有启发}

前面的分类来自自动驾驶，但本节要做迁移：机器人没有现成汽车市场，也没有天然车队数据。正因为没有这些前提，自动驾驶经验只能提供框架，不能直接提供答案。

具身智能还没有像汽车那样成熟的需求侧和车队规模。自动驾驶给机器人的启发不是照搬 robotaxi 或供应商模式，而是看清数据闭环必须和客户价值绑定。机器人若没有真实场景和出货，就没有持续数据；没有数据，就难以训练动作基础模型；没有模型能力，就无法提升客户价值。

\begin{warningbox}{自动驾驶经验不能机械复制}
汽车本身已有巨大市场，车能卖出去；通用机器人本体还没有同样成熟需求。机器人必须同时创造硬件需求、场景价值和数据闭环，难度比“把自动驾驶路线搬过来”更高。
\end{warningbox}

\subsection{本章小结}

自动驾驶商业模式的底层问题是：数据从哪里来，客户价值怎么形成，系统如何迭代。具身智能同样绕不开这些问题，只是它的需求、硬件和场景更分散。

